\documentclass{article}
\usepackage{graphicx}
\usepackage{amsmath}
\usepackage{amssymb}
\usepackage{amsfonts}
\usepackage{hyperref}
\usepackage{url}
\usepackage{listings}
\usepackage{xcolor}
\usepackage{float}

\lstset{
    language=Python,
    basicstyle=\ttfamily\small,
    keywordstyle=\color{blue},
    stringstyle=\color{red},
    commentstyle=\color{gray},
    showstringspaces=false,
    frame=single,
    breaklines=true
}



\title{Redes Neuronales Profundas (NN) - T01}
\author{ALEJANDRO ZARATE MACIAS}
\date{9 de Marzo 2026}

\begin{document}

\maketitle

\begin{abstract}
En esta tarea se abordan conceptos fundamentales de optimización numérica y álgebra lineal computacional como base para el estudio de las redes neuronales profundas. Se implementan y analizan distintos métodos de optimización, diferenciación automática y cálculo de autovalores, haciendo uso de PyTorch como herramienta principal.
\end{abstract}


% ========================================
% SECCIÓN 1
% ========================================
\section{Problema 1}

\subsection{Enunciado}

\section*{Problema 1}

Sea $\mathbf{x}$ y $\mathbf{1} := (1,1,\ldots,1)$ ambos en $\mathbb{R}^n$. Considere la función de Rosenbrock:

\[
f(\mathbf{x}) = \sum_{i=1}^{n-1} \left[ 100(x_{i+1} - x_i^2)^2 + (x_i - 1)^2 \right].
\tag{1}
\]

Calcule $\nabla f(\mathbf{1})$ utilizando diferenciación automática en modo reverso (a mano) para $n = 5$.

\subsection{Metodología}

\subsection{Resultados}

\subsection{Discusión}

\subsection{Conclusión}


% ========================================
% SECCIÓN 2
% ========================================
\section{Problema 2}

Haga un script en Python para calcular el autovalor de módulo máximo y su autovector correspondiente de una matriz simétrica $A \in \mathbb{R}^{n\times n}$. Use Pytorch en lugar de los paquetes Numpy/Scipy. Use el método de la potencia, descrito (por ejemplo) en \cite{golub1996matrix}, en la sección 8.2.1. A continuación, pruebe su programa con la siguiente matriz

\[
\left(
\begin{array}{rrrr}
-1.6407 & 1.0814 & 1.2014 & 1.1539 \\
1.0814 & 4.1573 & 7.4035 & -1.0463 \\
1.2014 & 7.4035 & 2.7890 & -1.5737 \\
1.1539 & -1.0463 & -1.5737 & 8.6944
\end{array}
\right).
\]

\subsection{Enunciado}

\subsection{Metodología}

\subsection{Resultados}

\subsection{Discusión}

\subsection{Conclusión}


% ========================================
% SECCIÓN 3
% ========================================
\section{Problema 3}

\subsection{Enunciado}

Haga un script en Python usando Pytorch para minimizar la función (1) para $n = 10$ utilizando el método de Newton. Grafique cómo el error disminuye a medida que el algoritmo itera.

\subsection{Metodología}

\subsection{Resultados}

\subsection{Discusión}

\subsection{Conclusión}


% ========================================
% SECCIÓN 4
% ========================================
\section{Problema 4}

\subsection{Enunciado}

Haga un script en Python usando Pytorch para minimizar la función (1) para $n = 10$ utilizando el método BFGS. Grafique cómo el error disminuye a medida que el algoritmo itera.

\subsection{Metodología}

\subsection{Resultados}

\subsection{Discusión}

\subsection{Conclusión}


% ========================================
% SECCIÓN 5
% ========================================
\section{Problema 5}

\subsection{Enunciado}

Haga un script en Python usando Pytorch para minimizar la función (1) para $n = 10$ utilizando el método de Adam. Grafque cómo el error disminuye a medida que el algoritmo iteira.

\subsection{Metodología}

\subsection{Resultados}

\subsection{Discusión}

\subsection{Conclusión}


% ========================================
% REFERENCIAS
% ========================================
\bibliographystyle{plain}
\bibliography{ref}

\end{document}
